\chapter{Introducere}
\section{Prezentarea temei}
Șahul este un veritabil joc de strategie care a antrenat mințile oamenilor din vremuri imemorabile și care încă dăinuie în ziua de astăzi. Fie că este jucat după regulile standard sau nu, acesta este destul de flexibil și ușor de înțeles încât a-i oferi jucătorului o experiență plăcută. Una dintre cele mai populare reguli este de a-i aloca fiecărui jucător un timp limitat în care aceștia pot acționa, depășirea acestuia cauzând sfârșitul partidei curente de șah în defavoarea celui rămas fără timp. Evident că se pune problema setării și gestionare a acestui timp printr-un dispozitiv, indiferent de natura acestuia (analog sau digital).

\section{Scop și obiective}
Pentru cel care vrea să își organizeze o partidă de șah acesta are la dispoziția sa o multitudine de opțiuni pe piață, dintre care vom aminti \textit{Digital Game Technology}. Digital Game Technology (DGT pe scurt) este olandeză ce produce ceasuri de șah aprobate de FIDE (Fédération Internationale des Échecs)) \cite{dgt_about}. Aceste ceasuri pot fi utilizate împreună cu table inteligente de șah sau separt. Mai nou acesta conțin un întreg sistem de calcul ce permit jucătorului să joace meciuri împotriva calculatorului fără a mai fi necesară conectarea la un PC \cite{dgt_computer}. \\
Scopul acestei lucrări este de a realiza interfațarea ceasului \textit{DGT 3000} cu un PC folosind un microcontroller de tip STM32F411.

\section{Cerințe}
\begin{enumerate}
	\item Soluția propusă va oferi o bibliotecă pentru interfața cu ceasurile de tip DGT 3000.
	\item Utilizatorul va folosi doar portul USB al calculatorului pentru a comunica cu ceasul de tip DGT 3000.
	\item Soluția propusă va oferi o aplicație pentru a facilita comunicația biderecțională prin portul USB dintre ceasul de tip DGT 3000 și calculatorul utilizatorului.
\end{enumerate}
